% ===== PRÉAMBULE CANONIQUE — à coller VERBATIM en tête de CHAQUE .tex =====
\documentclass[11pt,a4paper]{article}
\usepackage[utf8]{inputenc}\usepackage[T1]{fontenc}\usepackage{lmodern}
\usepackage[french]{babel}
\usepackage{amsmath,amssymb,mathtools}\usepackage{icomma}
\usepackage[version=4]{mhchem}          % \ce{...} : équations & formules
\usepackage{siunitx}\sisetup{locale=FR} % \qty{}{}, \si{}, \ang{}
\usepackage{chemfig}                    % \chemfig{...} : structures développées
\usepackage{xcolor}\usepackage{colortbl}
\usepackage{tikz}\usepackage{pgfplots}\pgfplotsset{compat=1.18}
\usepackage[most]{tcolorbox}\usepackage{enumitem}\usepackage{array,booktabs}
\usepackage[a4paper,margin=2.2cm]{geometry}
\usetikzlibrary{arrows.meta,calc,angles,quotes,positioning,patterns,backgrounds,shapes.geometric}
\definecolor{primary}{RGB}{222,49,99}\definecolor{primarydark}{RGB}{150,22,55}
\definecolor{defcol}{RGB}{0,114,178}\definecolor{methcol}{RGB}{52,92,148}
\definecolor{excol}{RGB}{0,135,90}\definecolor{attncol}{RGB}{200,40,55}
\tcbset{boxbase/.style={enhanced,breakable,boxrule=0.9pt,arc=2.5pt,
  left=3mm,right=3mm,top=2.3mm,bottom=2.3mm,fonttitle=\bfseries,coltitle=white}}
% --- boîtes TUTORIEL (noms EXACTS, ne pas renommer) ---
\newtcolorbox{cle}[1][]{boxbase,colback=defcol!6,colframe=defcol,
  title={\textsc{À retenir}\ifx\relax#1\relax\relax\else\textmd{ : \itshape #1}\fi}}
\newtcolorbox{methode}[1][]{boxbase,colback=methcol!7,colframe=methcol,sharp corners,
  title={\textsc{Méthode}\ifx\relax#1\relax\relax\else\textmd{ --- \itshape #1}\fi}}
\newtcolorbox{exemplebox}[1][]{boxbase,colback=excol!5,colframe=excol,
  title={\textsc{Exemple}\ifx\relax#1\relax\relax\else\textmd{ : \itshape #1}\fi}}
\newtcolorbox{piege}[1][]{boxbase,colback=attncol!6,colframe=attncol,
  title={\textsc{Piège}\ifx\relax#1\relax\relax\else\textmd{ : \itshape #1}\fi}}
\newtcolorbox{remarque}[1][]{boxbase,colback=black!4,colframe=black!45,
  title={\textsc{Remarque}\ifx\relax#1\relax\relax\else\textmd{ : \itshape #1}\fi}}
% --- boîtes EXERCICE / CORRIGÉ (noms EXACTS, auto-comptées) ---
\newtcolorbox[auto counter]{exo}[1][]{boxbase,colback=primary!4,colframe=primary,
  title={\textsc{Exercice}~\thetcbcounter\ifx\relax#1\relax\relax\else\textmd{ --- \itshape #1}\fi}}
\newtcolorbox[auto counter]{corrige}[1][]{boxbase,colback=excol!4,colframe=excol,
  title={\textsc{Corrigé}~\thetcbcounter\ifx\relax#1\relax\relax\else\textmd{ --- \itshape #1}\fi}}
\newcommand{\R}{\mathbb{R}}\newcommand{\N}{\mathbb{N}}\newcommand{\Z}{\mathbb{Z}}\newcommand{\ds}{\displaystyle}
% =========================================================================
\begin{document}
\sisetup{per-mode=symbol}
\begin{exo}[précipitation lors d'un mélange]
On mélange \qty{50}{\milli\litre} de nitrate de plomb \ce{Pb(NO3)2} à \qty{0.020}{\mole\per\litre} avec
\qty{50}{\milli\litre} d'iodure de potassium \ce{KI} à \qty{0.020}{\mole\per\litre}. On donne
$K_{\mathrm{ps}}(\ce{PbI2})=\num{7.1e-9}$.
\begin{enumerate}[label=\alph*)]
  \item Calcule $[\ce{Pb^{2+}}]$ et $[\ce{I-}]$ juste après mélange (pense à la dilution).
  \item Calcule $Q$ et compare à $K_{\mathrm{ps}}$ : \ce{PbI2} précipite-t-il ?
\end{enumerate}
\end{exo}

\begin{exo}[effet d'ion commun sur \ce{AgCl}]
Une solution saturée de \ce{AgCl} ($K_{\mathrm{ps}}=\num{1.8e-10}$) reçoit des ions chlorure jusqu'à
$[\ce{Cl-}]=\qty{0.10}{\mole\per\litre}$ (à l'équilibre).
\begin{enumerate}[label=\alph*)]
  \item Rappelle la solubilité $s_0$ dans l'eau pure.
  \item Calcule la nouvelle solubilité $s'=[\ce{Ag+}]$.
  \item Par quel facteur la solubilité a-t-elle été divisée ?
\end{enumerate}
\end{exo}

\begin{exo}[ion commun avec un anion doublement présent]
On dissout \ce{PbBr2} ($K_{\mathrm{ps}}=\num{9.2e-6}$) dans une solution de \ce{KBr} telle que
$[\ce{Br-}]=\qty{0.50}{\mole\per\litre}$.
\begin{enumerate}[label=\alph*)]
  \item Écris $K_{\mathrm{ps}}$ et isole la solubilité $s'$ (attention à l'exposant de $[\ce{Br-}]$).
  \item Calcule $s'$.
\end{enumerate}
\end{exo}

\begin{exo}[concentration seuil de précipitation]
Une solution contient $[\ce{Ag+}]=\qty{1.0e-3}{\mole\per\litre}$. On y ajoute progressivement des ions
\ce{Cl-}. On donne $K_{\mathrm{ps}}(\ce{AgCl})=\num{1.8e-10}$.
\begin{enumerate}[label=\alph*)]
  \item À partir de quelle concentration $[\ce{Cl-}]$ la précipitation de \ce{AgCl} commence-t-elle ?
  \item Que se passe-t-il en dessous de cette valeur ?
\end{enumerate}
\end{exo}

\begin{exo}[y a-t-il précipitation de \ce{CaCO3} ?]
On donne $K_{\mathrm{ps}}(\ce{CaCO3})=\num{3.3e-9}$. Pour chaque mélange, calcule $Q$ et conclus.
\begin{enumerate}[label=\alph*)]
  \item $[\ce{Ca^{2+}}]=\qty{1.0e-3}{\mole\per\litre}$, $[\ce{CO3^{2-}}]=\qty{1.0e-5}{\mole\per\litre}$.
  \item $[\ce{Ca^{2+}}]=\qty{1.0e-4}{\mole\per\litre}$, $[\ce{CO3^{2-}}]=\qty{1.0e-5}{\mole\per\litre}$.
\end{enumerate}
\end{exo}
\end{document}
